\section*{Цели и задачи работы}

\textit{Цели работы:} измерение сопротивлений объемных и тонкопленочных резисторов; исследование зависимостей удельных электрических сопротивлений и их температурных коэффициентов от температуры и состава резистивных материалов; исследование зависимостей термоЭДС термопар от разности температур контактов.

\textit{Задачи работы:}
\begin{enumerate}
    \item Измерить сопротивление пленочных и проволочных резисторов при комнатной температуре.
    \item Зафиксировать геометрические размеры образцов, указанные на лицевой панели стенда.
    \item Провести измерения сопротивлений резисторов из никеля, меди и константана в диапазоне температур от комнатной до 150 °С.
    \item Провести измерения термоЭДС термопар (медь-железо, медь-константан, медь-манганин) в том же диапазоне температур.
    \item Занести все измеренные значения в таблицы протокола.
\end{enumerate}

\section*{Краткие теоретические сведения}

К основным электрическим характеристикам проводниковых материалов относят удельное сопротивление $\rho$ и температурный коэффициент удельного сопротивления $\alpha_\rho$. Удельное сопротивление в рамках классической электронной теории проводимости металлов определяется выражением:
\begin{equation}
    \rho = \frac{m\bar{u}}{e^2 n_0 \bar{\lambda}},
    \label{eq:rho_def}
\end{equation}
где $m$ и $e$ — масса и заряд электрона, $n_0$ — концентрация свободных электронов, $\bar{u}$ — средняя скорость их теплового движения, $\bar{\lambda}$ — средняя длина свободного пробега.

Относительное изменение удельного сопротивления при изменении температуры на один кельвин называют температурным коэффициентом удельного сопротивления (ТКС):
\begin{equation}
    \alpha_\rho = \frac{1}{\rho}\frac{d\rho}{dT}.
    \label{eq:tcs_def}
\end{equation}

Для тонких пленок вводят понятие удельного поверхностного сопротивления (сопротивление квадрата поверхности):
\begin{equation}
    R_\square = \rho / d,
    \label{eq:rho_surf}
\end{equation}
где $d$ — толщина пленки.

При контакте двух различных проводников возникает термоэлектродвижущая сила (термоЭДС), которая в небольшом температурном интервале пропорциональна разности температур контактов (спаев):
\begin{equation}
    \Delta U = a(t_2 - t_1),
    \label{eq:thermo_emf}
\end{equation}
где $a$ — удельная термоэлектродвижущая сила.

\section*{Порядок выполнения работы}

\subsection*{1. Определение удельного сопротивления при комнатной температуре}
\begin{enumerate}
    \item \textit{Подготовка измерительного стенда.}
    \begin{itemize}
        \item \textit{Что сделать:} Включить мультиметр в сеть.
        \item \textit{Как сделать:} Перевести тумблер питания на задней панели прибора в положение «Вкл».
        \item \textit{Нюансы:} Дать прибору прогреться в течение 5 минут для стабилизации показаний. Тумблер «Установка температуры» на термостате должен быть выключен.
    \end{itemize}
    \item \textit{Настройка мультиметра.}
    \begin{itemize}
        \item \textit{Что сделать:} Установить режим измерения сопротивления.
        \item \textit{Как сделать:} Переключатель режимов мультиметра установить в положение «Ом», переключатель пределов измерений — в положение «200 Ом».
    \end{itemize}
    \item \textit{Измерение сопротивлений и фиксация размеров.}
    \begin{itemize}
        \item \textit{Что сделать:} Последовательно измерить сопротивление пленочных (поз. 1...3) и проволочных (поз. 6...10) резисторов.
        \item \textit{Как сделать:} Для измерения сопротивления конкретного образца необходимо нажать и удерживать соответствующую ему кнопку на лицевой панели стенда.
        \item \textit{Нюансы:} В случае перегрузки прибора (индикация «1» на левом краю дисплея) следует переключить мультиметр на более высокий предел измерения.
        \item \textit{Что сделать:} Занести результаты измерений сопротивлений и геометрические размеры образцов (указаны на панели) в таблицы \cref{p:pr13} и \cref{p:pr610} протокола.
    \end{itemize}
\end{enumerate}

\subsection*{2. Определение температурных зависимостей сопротивления и термоЭДС}
\begin{enumerate}
    \item \textit{Измерение при комнатной температуре.}
    \begin{itemize}
        \item \textit{Что сделать:} Измерить начальные значения сопротивлений резисторов $R_1$ (Никель), $R_2$ (Медь), $R_3$ (Константан) и термоЭДС всех термопар.
        \item \textit{Как сделать:} Провести измерения, как описано в предыдущем пункте. Для измерения термоЭДС переключить мультиметр в режим измерения постоянного напряжения (положение «V=»).
        \item \textit{Нюансы:} Записать полученные значения как первую точку (при комнатной температуре) в таблицы \cref{p:rt} и \cref{p:eds}.
    \end{itemize}
    \item \textit{Нагрев образцов.}
    \begin{itemize}
        \item \textit{Что сделать:} Включить нагрев термостата и установить минимальную температуру.
        \item \textit{Как сделать:} Включить тумблер «Установка температуры» и установить регулятор температуры в положение «min».
    \end{itemize}
    \item \textit{Проведение измерений в цикле нагрева.}
    \begin{itemize}
        \item \textit{Что сделать:} Проводить измерения сопротивлений и термоЭДС с шагом aproximadamente 10-15 °С вплоть до 150 °С.
        \item \textit{Как сделать:} После стабилизации температуры на каждом шаге (контролируется по показаниям термометра) провести измерения сопротивлений $R_1, R_2, R_3$ и термоЭДС трех термопар.
        \item \textit{Нюансы:} При каждом измерении термоЭДС необходимо регистрировать температуру не только горячего спая (в термостате), но и холодного спая (вне термостата).
        \item \textit{Что сделать:} Все результаты измерений заносить в таблицы \cref{p:rt} и \cref{p:eds}.
    \end{itemize}
    \item \textit{Завершение работы.}
    \begin{itemize}
        \item \textit{Что сделать:} Завершить измерения и выключить оборудование.
        \item \textit{Как сделать:} По достижении максимальной температуры регулятор температуры поставить в крайнее левое положение, выключить тумблер «Установка температуры» и выключить мультиметр.
    \end{itemize}
\end{enumerate}

% --- Список иллюстраций ---
% fig:front_panel: Схематичное изображение лицевой панели установки.
% fig:circuit: Принципиальная схема подключения образцов к мультиметру.